\documentclass[10pt,letterpaper,final]{report}
\usepackage[utf8]{inputenc}
\usepackage{amsmath}
\usepackage{amsfonts}
\usepackage{amssymb}
\usepackage{amsthm}
\renewcommand\qedsymbol{$\blacksquare$}
\usepackage{enumerate}
\usepackage{hyperref}
\usepackage{pdfpages}
\usepackage{graphics}
\usepackage{graphicx}
\usepackage{tikz}
\usepackage{tikz-qtree}
\usetikzlibrary{automata,arrows}
\usepackage{mdframed}
\usepackage[
  backend=biber,
  style=authoryear,
  sorting=nyt          % sort by name, year, title
]{biblatex}
\addbibresource{references.bib}

\usepackage{minted}
\usemintedstyle{colorful}

% Based on template from COSC-417 at Towson University, originally authored by Marius Zimand

\begin{document}

\fbox{
  \vbox{
    \begin{flushleft}
      Jonathan Llovet \\  % authors' names
      EN.605.202.81 \\  %class
      2026-06-09 \\  % date
      Module 3 - Homework 3 \\ % ID
    \end{flushleft}
    \center{\Large{\textbf{Data Structures - Recursion}}}
    %\end{mdframed}
  } % end vbox
} % end fbox
\vline

\noindent\textbf{Problem 1:}

Write a recursive algorithm to compute \verb|a+b|, where \verb|a| and \verb|b| are nonnegative integers.

\noindent\textbf{Answer:}

\begin{mdframed}
\begin{verbatim}
function add(a,b) {
  if b == 0 {
    return a
  } else {
    return add(a+1,b-1)
  }
}
\end{verbatim}
\end{mdframed}

\pagebreak
\noindent\textbf{Problem 2:}

Let \verb|A| be an array of integers. Write a recursive algorithm to compute the average of the elements of the array. Solutions calculating the sum recursively, instead of the average, are worth fewer points.

\noindent\textbf{Answer:}

The algorithm below takes advantage of the identity:
\[
\frac{a + b + c}{n} = \frac{a}{n} + \frac{b}{n} + \frac{c}{n}
\]

In the function below, I mean \verb|tail(a)| as a function that returns an array \verb|a'| that is all of the elements of \verb|a| from index 1 to the end. For example: \[\verb|tail([1,2,3]) -> [2,3]|.\]

\begin{mdframed}
\begin{verbatim}
function running_average(A){
    // Length of original A used for running average
    len_A = A.length
    // Initialize the running average to 0
    running = 0
    // Define internal procedure for perfoming the recursion
    function avg(A, running, len_A){
        // Base case
        if A.length == 0 {
            return running
        }
        // Recursive case
        running = running + A[0]/len_A
        return avg(tail(A), running, len_A)
    }
    // Call it
    return avg(A, running, len_A)    
}
\end{verbatim}
\end{mdframed}

Or, equivalently and more simply:

\begin{mdframed}
\begin{verbatim}
function avg(A, running, len_A){
    if A.length == 0 {
        return running
    }
    // length of original A passed down
    running = running + A[0]/len_A
    return avg(tail(A), running, len_A)
}
\end{verbatim}
\end{mdframed}

This can be called as: \verb|avg(A, 0, A.length)|


\pagebreak
\noindent\textbf{Problem 3:}

If an array contains \verb|n| elements, what is the maximum number of recursive calls made by the binary search algorithm?

\noindent\textbf{Answer:}

This is discussed in our textbook \cite{zybooks-ds2026} in section 3.6 ``Analyzing the time complexity of recursive algorithms''. What is the reasoning? For an input of size \verb|n|, binary search can be analyzed into two cases, the base case and the recursive case. In the base case, the position is found or is determined to be absent from the array. That requires a single call to the binary search function. In the recursive case, binary search is called on half of the list in the current scope, i.e., $n/{2}^1$, $n/{2}^2$, $n/{2}^3$, and so forth. The number of these recursive calls is dependent on the size of $n$. Since the size of the sublist that is searched is halved each time that the binary search function is called, $\log_{2}n$ calls are required for the recursive case. Therefore, presuming a correct implementation and that the input list is sorted, for a sorted input list of length $n$, there are $\lfloor 1 + \log_{2}n \rfloor$ calls to binary search required in the worst case to find the item or to determine that it is not in the list. When the input list has a size that is not a power of $2$, the $\log_{2}n \not \in \mathbb{Z}$. But the number of steps taken by the algorithm must be some $x \in \mathbb{Z}^+$. In light of this, we round down here to the highest integer less than the actual value of $1 + \log_{2}n$. The script below can be used to compare values of $1 + \log_{2}n$ for various values of $n$ to validate this empirically.

\inputminted[breaklines=true,fontsize=\footnotesize]{python}{m3p3-bs.py}

\pagebreak

A sample of the output from the the script above for $1 \le x \le 10$:

\inputminted[breaklines=true,fontsize=\footnotesize]{json}{experiment-results.json}
\nocite{OverleafcomCodeHighlightingWithMinted}


\pagebreak
\noindent\textbf{Problem 4:}

The expression \verb|m % n| yields the remainder of \verb|m| upon (integer) division by \verb|n|. Define the greatest common divisor (GCD) of two integers \verb|x| and \verb|y| by:
\begin{itemize}
  \item $gcd(x, y) = y$ if $(y \le x$ and $x\%y == 0)$
  \item $gcd(x, y) = gcd(y, x)$ if $(x < y)$
  \item $gcd(x, y) = gcd(y, x\%y)$ otherwise
\end{itemize}

Write a recursive method to compute \verb|gcd(x,y)|.

\noindent\textbf{Answer:}

The function is very straightforward from this specification:

\begin{mdframed}
\begin{verbatim}
function gcd(x, y){
  if y <= x and x % y == 0 {
    return y
  }
  if x < y {
    return gcd(y,x)
  }  
  else {
    return gcd(y, x%y)
  }
}
\end{verbatim}
\end{mdframed}

\pagebreak
\noindent\textbf{Problem 5:}

A generalized Fibonacci function is like the standard Fibonacci function, except that the starting points are passed in as parameters. Define the generalized Fibonacci sequence of \verb|f0| and \verb|f1| as the sequence

\verb|gfib( f0, f1, 0), gfib(f0, f1, 1), gfib(f0, f1, 2), ...,| where
\begin{itemize}
  \item \verb|gfib(f0, f1, 0) = f0|
  \item \verb|gfib(f0, f1, 1) = f1|
  \item \verb|gfib(f0, f1, n) = gfib(f0, f1, n-1) + gfib(f0, f1, n-2) if n > 1|
\end{itemize}

Write a recursive method to compute \verb|gfib(f0,f1,n)|.

\noindent\textbf{Answer:}

This problem is similarly straightforward from the specification.

\begin{mdframed}
\begin{verbatim}
function gfib(f0, f1, n) {
  if n == 0 {
    return f0
  }   
  if n == 1 {
    return f1
  }   
  if n > 1 {
    return gfib(f0, f1, n-1) + gfib(f0, f1, n-2)
  }       
}
\end{verbatim}
\end{mdframed}

\pagebreak
\noindent\textbf{Problem 6:}

Ackerman's function is defined recursively on the nonnegative integers as follows:
\begin{itemize}
  \item \verb|a(m, n) = n + 1 if m = 0|
  \item \verb|a(m, n) = a(m-1, 1) if m \ne 0, n = 0|
  \item \verb|a(m, n) = a(m-1, a(m, n-1)) if m \ne 0, n \ne 0|
\end{itemize}
Using the above definition, show that \verb|a(2,2)| equals \verb|7|.

\noindent\textbf{Answer:}

I wrote the script below to perform the calculations for Ackerman's function and to trace the depth of the recursive function calls.

\inputminted[breaklines=true,fontsize=\footnotesize]{python}{m3p6-ackerman.py}

\pagebreak

\noindent My program for tracking Ackerman's function's execution produces this trace:
\inputminted[breaklines=true,fontsize=\tiny]{text}{ackerman-2-2.txt}


\pagebreak
\noindent\textbf{Problem 7:}

Convert the following recursive program scheme into an iterative version that does not use a stack. \verb|f(n)| is a method that returns \verb|TRUE| or \verb|FALSE| based on the value of \verb|n|, and \verb|g(n)| is a method that returns a value of the same type as \verb|n| (without modifying \verb|n| itself).

\begin{mdframed}
\begin{verbatim}
int rec(int n) {
  if (f(n) == FALSE) {
    /* any group of statements that
    do not change the value of n */
    return (rec(g(n)));
  }//end if
  return (0);
}//end rec
\end{verbatim}
\end{mdframed}

\noindent\textbf{Answer:}

The recursive function is called with g(n) each time the function is called after the initial invocation. Presuming that \verb|f(n)| returns \verb|FALSE| so that \verb|rec(g(n))| is called, if \verb|g(n)| returns a value equal to \verb|n| every time, then the function \verb|rec| will continue to recurse forever or until the resources of the machine are exhaused. If however \verb|g(n)| returns a different value than \verb|n|, then the recursion may eventually terminate, depending on whether the value of \verb|g| composed with itself repeatedly approaches a value \verb|x| such that \verb|f(x)| returns \verb|TRUE| or not. This can happen, for example, in ways that we have seen in the other exercises, where there is a recurrence relation in which \verb|g| composed with itself several times causes its argument to decrement towards \verb|1| or \verb|0|. Ultimately, in this general form, the termination of the function overall is dependent on the value of \verb|f(n)|. An iterative scheme that achieves this behavior is a while loop that has \verb|f(n)| as the predicate for the while loop to check, then \verb|g(n)| as a function that may increment or modify the value of \verb|n| that is in scope for \verb|f(n)|. The same characteristics about the possible termination or infinite looping of the outer function hold for the iterative function below as they did for the recursive case in the prompt.


\begin{mdframed}
\begin{verbatim}
int iterative(int n){
  while (f(n) == FALSE) {
    /*actions that don't mutate n*/
    n = g(n)
  }
  return 0
}
\end{verbatim}
\end{mdframed}


\pagebreak

\section*{Source Code}
The full source code, including LaTeX, tooling, other artifacts, and git history is available on my Github repository for the class here:
\begin{center}
  \url{https://github.com/jllovet/jhu-en-605-202-su26-data-structures}
\end{center}

\printbibliography

\end{document}